\section{Emergence}

Everything above the tones is a pattern, and the question this section answers is when a pattern is real. A breaking wave is real though it is only moving water, and a traffic jam is real though it is only cars, because each has a life of its own, a behavior you can predict at its own level without tracking every piece beneath. The base of this theory is tones and one law. Everything else, an electron, a cell, a self, a mind, is a pattern in them, and the work of the upper half of the paper is to say which patterns earn the name and how they come to be.

Three marks separate a real pattern from a passing accident, and all three are measured rather than asserted. A real pattern has causal power at its own level, so its coarse description forecasts its future as well as the fine one does, and sometimes better. It persists, holding its shape while the tones beneath it turn over, the way a flame keeps its form on always-new gas. And it is integrated, more than the sum of its parts, so that cutting it apart destroys something the parts alone do not carry. A clustering of tones with none of these is weather, not a thing. A pattern with all three is an object at its level, as solid there as an atom is at its own.

Patterns arise by a small set of routes, and naming them is most of the science. The plainest is a traveling shape the law re-creates step after step, the glider already met, the seed of every particle. Above it is topological protection, a knot or a twist in the tones that no local move can undo, so it persists for the same reason a knot in a rope survives shaking. Above that is coarse-graining, the averaging of many tones into one effective quantity, which is how a smooth field arises from whole-number tones and how each higher level is born from the one below. The discipline that governs all of it is the rule of levels. The base is tones, ordinary physics is their coarse-grained face at the cusp, life and mind higher still, and a phenomenon must be read at its own rung. A clump of raw tones is not gravity, and a single high tone is not a particle, however much each looks the part. Skipping a level is the commonest way to fool yourself, and the paper works never to.

\begin{result}[When a pattern is real]
A pattern counts as a thing at its level when it has causal power there, persists while its substrate turns over, and is irreducible to its parts. Each is measured, with a control that fails, so the line between a real emergent object and a suggestive clustering is drawn by measurement and not by eye. \cite{pollard2026vibetest}
\end{result}

The state of the program is that the lowest emergent levels are in hand and the climb past them is open. The base on its own churns, and a lasting self does not appear from the pure law without the further ingredients the next sections supply, a result kept in plain view rather than wished away. One natural way to coarse-grain the curved bulk fails to produce the tower of scales that real physics shows, and that failure is reported, not buried.

There is a clean reason that tower must live above the base and not within it. The base law is a bijection, exactly reversible, and a bijection has no slack, no two fine states ever falling together into one coarse state. So at the spotless base no coarse level can carry more causal power than the fine one beneath it, and genuine emergence, the kind where a higher law steers the lower, needs a lossy step, a little forgetting, before it can appear at all. That forgetting is what the bath supplies. So the hierarchy of selves is an achievement of the effective layer, real and measured once a little is allowed to be lost, and absent from the bare reversible law, which is also why the attempt to read a tower straight off the reversible bulk by averaging in shells returned nothing the plain spreading of heat did not already give, and was built, measured, and withdrawn.

So the forward work is clear, and it is mostly simulation. Find the coarse-graining of the four-dimensional bulk that yields a clean layer tower. Settle whether a self, once formed, radiates or binds, by running it longer and larger than has yet been done. Pin the second conserved quantity a fully relativistic massless mode needs. And push the substrate to the larger sizes where an emergent law either holds or breaks, since every claim here is the kind a bigger computation can confirm or refute. The theory is built to be run, and the next results will come from running it.
